\section{Tests}

The tests carried out during development began with general electrical verification of the assembled PCB in order to ensure correct soldering and reliable basic operation. At an early stage, a multimeter was used to check the main power path across the board, starting from the 7.4 V battery, through the H-bridge and the first voltage regulator, which generated the 5 V supply. This 5 V line was then verified as the supply for the three analog infrared sensors and as the input to a second regulator providing 3.3 V for the remaining digital components, including the microcontroller. After these static supply checks, the verification was extended by oscilloscope measurements in order to assess dynamic electrical signals at relevant module interfaces, for example at PWM outputs and other time-varying signals. In this way, the initial test phase did not only confirm the presence of the required supply voltages, but also provided a general board-level check that the major components received the expected electrical signals before more specific tests were carried out.

The purpose of this chapter is to describe how the individual subsystems were validated during development. For this reason, the emphasis is on the test procedures, debugging steps, and validation scenarios used during commissioning. The final technical outcomes and the selected operating parameters are summarized separately in the following results chapter.

\subsection{Sensor Tests}

The sensor tests were comparatively simple, since the infrared sensors functioned reliably from the beginning and therefore did not require extensive dedicated tuning. After the sensors had been connected to the system, all three raw sensor signals were first observed via UART output in order to verify that valid readings were received and updated correctly during operation. The initial goal of these tests was therefore not detailed calibration, but confirmation that the sensing chain from analogue input to software interpretation worked as intended.

To evaluate the sensors under realistic operating conditions, the robot was then placed inside the maze and positioned close to surrounding walls. These measurements were used to compare free-space behaviour with the more relevant operating range inside the maze and to determine whether the sensors provided a sufficiently clear distinction between wall and no-wall conditions for the intended navigation task. Since the software relies primarily on threshold-based wall detection rather than on precise metric distance reconstruction, this practical distinction was more important than an exact calibration curve.

For easier debugging during deployment, the on-board LEDs were also used as direct visual indicators of wall detection, so that the sensor behavior could be checked without relying on a UART connection. In this configuration, the blue LED indicated detection by the left sensor, the red LED indicated detection by the right sensor, and the green LED indicated detection by the front sensor. This made it possible to verify the switching behavior of the wall-detection logic directly on the robot during operation in the maze.

Overall, the sensor system required only limited dedicated testing, since all three sensors worked essentially out of the box. The main outcome of the sensor tests was therefore not a redesign of the sensing concept, but the validation of the existing sensor arrangement and the derivation of a practical detection threshold for later use in the control and navigation logic.

\subsection{Controller Tests}

The controller tests were carried out incrementally, beginning with the inner wheel-speed control and later extending to the wall-based trim correction used during straight driving. As a prerequisite, the motor actuation had already been implemented in normalized form, such that motor commands could be specified in the range from $-1$ to $1$ and were automatically limited to the admissible actuation range. Building on this interface, the first controller-related tests focused on verifying that a commanded translational speed in $\mathrm{mm/s}$, provided through the function \texttt{setDriveSpeedMmps()}, could be mapped to a corresponding fixed motor actuation and that the drivetrain responded accordingly. At this stage, the purpose was not yet closed-loop regulation, but rather to confirm that the command path from software input to physical wheel motion worked as intended.

For observation during testing, a verbose UART output was available from the beginning and was transmitted every $10\,\mathrm{ms}$ from the timer-interrupt-based controller update. This output included the measured left and right wheel velocities, the sensor readings, and further internal variables relevant for debugging, so that no additional instrumentation had to be added during tuning. Once the basic control path had been verified, the structure was extended to a closed-loop PI controller for wheel-speed regulation inside \texttt{updateWheelSpeedController()}, which was called periodically from \texttt{updateController()}. The tuning process then began with both controller gains set initially to zero. First, a proportional sweep was carried out on the right wheel with $K_I = 0$ and a simple two-step speed command, in order to identify the range in which visible speed regulation began and the point at which stuttering became excessive. After a suitable proportional range had been identified, the integral gain was increased incrementally while a wider cyclic speed sequence was applied. This made it possible to compare the reduction of steady-state deviation against the onset of increasingly aggressive behaviour. The left wheel was then tuned starting from the parameter region identified on the right side in order to check whether a common parameter set could be used. A summary of this empirical tuning sequence and the associated qualitative observations is provided in Appendix~A.

In principle, separate parameter sets for both wheels would be justified, since the two motors and the associated drivetrain are not perfectly identical. For this reason, both wheels were initially considered independently during tuning. The observations summarized in Appendix~A show, however, that both sides behaved similarly enough around the selected tuning region that a common parameter set could be used in the final implementation. This simplified the controller structure while still allowing satisfactory straight driving within the scope of the project.

After the single-wheel tests, both controlled wheels were evaluated together on different surfaces, including the floor, the maze surface, and a desk. The purpose of these tests was to determine whether the robot still exhibited sufficiently straight motion under changing contact conditions. Since the drive behavior remained acceptable on all tested surfaces, this part did not require extensive further investigation.

Once the inner PI controller had reached a stable and usable state, testing was extended to the wall-based trim correction. A first proof-of-concept test was carried out in open space, where no surrounding wall was present and both side sensors therefore produced approximately zero readings. Under these conditions, no trim correction should be generated. To introduce a controlled lateral disturbance, a wooden slab was manually placed near either the left or the right side sensor while the robot was driving. The expected behavior was that the robot would steer away from the detected obstacle. During this test, it was observed that the correction direction was initially inverted, such that the robot turned toward the slab instead. This revealed a sign error in the correction logic, which was then reversed.

After this initial validation, the trim-based correction was tested under more realistic maze conditions. For this purpose, the robot was placed in the maze and repeatedly driven along a hardcoded square path. The test logic was implemented such that the robot drove straight until a front wall was detected, then stopped the straight-drive controller, executed an approximately $90^\circ$ left turn using \texttt{turnLeft90()}, and continued driving. This setup proved useful because the turning motion was only approximate, so that the robot did not always enter the next corridor in perfect alignment. As a result, the subsequent straight segment required active correction in order to return the robot toward a stable path within the corridor. The square path therefore provided a repeatable test scenario for evaluating whether the correction logic could compensate such misalignment robustly.

Initially, the trim correction inside \texttt{updateWheelSpeedController()} had been implemented as a PI controller. However, no satisfactory parameter set could be found that produced both fast and stable correction behavior. In practice, the robot exhibited pronounced lateral oscillations and followed a snake-like path along the corridor. After reevaluation in discussion with Dr.\ Alexander Lenz, the trim controller was reduced to a proportional controller. This modification led to significantly improved behavior and allowed the correction to react more directly to lateral deviations without introducing the previously observed oscillatory motion.

In addition, the correction logic had to account for the fact that not all maze sections provide the same lateral wall structure, so trim-based corrections were applied only when a valid side-wall reference was available. To improve practical robustness, filtering and slew-rate limitation were also introduced during tuning in order to reduce noise sensitivity and abrupt actuator changes.

Overall, controller testing was considered successful once the wheel-speed controller tracked commanded velocities reliably, the robot maintained sufficiently straight motion in practice, and the wall-based outer correction enabled repeated square-path execution in the maze without substantial loss of alignment. These tests established a sufficiently stable basis for later sensor-related testing and further system integration.

\subsection{Navigation Tests}

The navigation behaviour was validated experimentally in several incremental stages. Since no dedicated automated test framework was implemented in the code, verification was mainly carried out through practical trials with the physical robot. At the beginning, only a very small maze and simple controller settings were used in order to reduce complexity and make individual effects easier to observe. This made it possible to focus first on the relation between sensor feedback, travelled distance, and the internal cell-based position estimate.

An important early step was to drive the robot at very low speed and stop it whenever the software detected that a cell center had been reached. At the same time, UART output was used to indicate when such a detection occurred. This allowed a direct comparison between the internal software decision and the real physical position of the robot in the maze.

In addition to odometry-based checks, the front sensor values were examined while the robot was positioned at cell centers. Based on these measurements, a suitable value range was defined in which the sensor feedback could be interpreted as a valid indication of a cell center. This helped to make the detection more robust and reduced the risk of orientation errors during exploration. After these basic mechanisms worked reliably, the overall difficulty of the test scenarios was increased step by step.

The exploration logic was also introduced incrementally. At first, the robot was only able to explore direct neighboring cells, without any recursive path calculation through the known maze. This simplified the behavior and made it easier to isolate problems in localisation, queue handling, and wall detection. Only after these subsystems appeared sufficiently robust, the recursive search was implemented, allowing the robot to calculate a path to queued cells that were no longer directly adjacent to its current position.

Once the full exploration procedure was functional, the remaining work focused on fine tuning. One observed issue was stuttering caused by the controller, which introduced noise into the distance measurement and therefore affected the reliability of cell-center detection. A misestimation of the cell center occurred particularly during longer driving periods, as the accumulated error increased over time. This required further adjustment of the controller behaviour and the related thresholds. After that, the robot speed was increased and the previous stop-and-check strategy was removed, so that planning and decision making could take place while the robot was still driving. This represented a more realistic operating mode and better matched the intended final behaviour of the system.

Finally, different maze layouts were tested in order to verify that the robot was able not only to explore reliably, but also to use the generated map correctly for the final run. In particular, these tests were used to check whether the robot consistently selected the optimal route to the predefined goal after exploration had finished.
